\subsection{Выдвижение гипотез}

Рассмотрим выборку \( \vec{X}=(X_1,\ldots,X_n)^T\sim\mathcal{F} \), где \( n = 100 \):

\begin{table}[h!]
\centering
\begin{tabular}{|c|c||c|c||c|c||c|c||c|c|}
\hline
№ & Значение & № & Значение & № & Значение & № & Значение & № & Значение \\ \hline \hline
1  & 0.32 & 21 & 0.23 & 41 & 0.40 & 61 & 0.35 & 81 & 0.54 \\ \hline
2  & 0.02 & 22 & 1.57 & 42 & 0.61 & 62 & 0.09 & 82 & 1.65 \\ \hline
3  & 0.14 & 23 & 0.06 & 43 & 0.11 & 63 & 0.16 & 83 & 1.63 \\ \hline
4  & 0.52 & 24 & 1.10 & 44 & 0.06 & 64 & 0.27 & 84 & 0.23 \\ \hline
5  & 0.14 & 25 & 0.39 & 45 & 0.13 & 65 & 0.41 & 85 & 0.67 \\ \hline
6  & 1.62 & 26 & 0.49 & 46 & 0.07 & 66 & 0.21 & 86 & 0.25 \\ \hline
7  & 0.28 & 27 & 0.34 & 47 & 0.18 & 67 & 0.37 & 87 & 0.00 \\ \hline
8  & 0.22 & 28 & 1.48 & 48 & 0.03 & 68 & 0.57 & 88 & 0.23 \\ \hline
9  & 0.07 & 29 & 0.06 & 49 & 0.66 & 69 & 0.31 & 89 & 0.43 \\ \hline
10 & 0.73 & 30 & 0.52 & 50 & 2.23 & 70 & 0.39 & 90 & 0.01 \\ \hline
11 & 0.83 & 31 & 0.31 & 51 & 0.04 & 71 & 0.30 & 91 & 0.10 \\ \hline
12 & 0.26 & 32 & 0.37 & 52 & 0.69 & 72 & 0.23 & 92 & 0.12 \\ \hline
13 & 1.09 & 33 & 0.15 & 53 & 0.06 & 73 & 0.56 & 93 & 2.04 \\ \hline
14 & 0.22 & 34 & 0.36 & 54 & 0.03 & 74 & 0.08 & 94 & 0.19 \\ \hline
15 & 1.43 & 35 & 0.45 & 55 & 0.96 & 75 & 0.14 & 95 & 0.07 \\ \hline
16 & 0.26 & 36 & 0.89 & 56 & 0.73 & 76 & 0.02 & 96 & 1.04 \\ \hline
17 & 0.25 & 37 & 0.35 & 57 & 0.67 & 77 & 0.58 & 97 & 0.17 \\ \hline
18 & 0.23 & 38 & 0.45 & 58 & 0.56 & 78 & 0.95 & 98 & 0.21 \\ \hline
19 & 0.05 & 39 & 0.15 & 59 & 0.02 & 79 & 0.49 & 99 & 0.55 \\ \hline
20 & 0.55 & 40 & 0.27 & 60 & 0.52 & 80 & 0.00 & 100& 0.97 \\ \hline
\end{tabular}
\caption{Выборка некоторой случайной величины \( X \)}
\end{table}

Нетрудно заметить, что значения случайной величины составляют некоторое подмножество \( \mathbb{R} \). Из этого делаем вывод, что \( X \) подчиняется некоторому \textit{непрерывному закону распределения} \( \mathcal{F} \).

\textbf{Гипотеза} \( H \) --- это любое предположение о распределении наблюдений. Выдвинем следующие гипотезы:
\[
    \begin{aligned}
        &H_1=\{ \mathcal{F}= \textrm{Exp}(\lambda) \}\\
        &H_2=\{ \mathcal{F}\neq \textrm{Exp}(\lambda) \}
    \end{aligned},
\]

где:
\begin{itemize}
    \item \( \lambda=1.41 \) --- параметр экспоненциального распределения;
    \item \( H_1 \) --- основная \textit{(простая)} гипотеза;
    \item \( H_2 \) --- альтернативная \textit{(сложная)} гипотеза.
\end{itemize}